\documentclass[10pt,a4paper]{amsart}
\usepackage[margin=19mm]{geometry}
\usepackage{amsmath,amssymb,mathtools}
\usepackage[scr=rsfs,cal=euler]{mathalfa}
\usepackage{xfrac}
\usepackage[unicode,hidelinks,bookmarks=false]{hyperref}
\newcommand{\ie}{i.e.\ }
\newcommand{\cf}{cf.\ }
\newcommand{\resp}{respectively\ }
\newcommand{\Subsection}{\subsection}
\providecommand{\nobreakdash}{\nobreak\hbox{-}}
\setlength{\emergencystretch}{2em}
\multlinegap0pt
\usepackage{bm}
\usepackage{bbm}
\usepackage{dsfont}
\usepackage{xspace}
\usepackage{cancel}
\usepackage{mathtools}
\usepackage[shortlabels]{enumitem}
\usepackage{amsthm}
\makeatletter
\@ifundefined{theorem}{\newtheorem{theorem}{Theorem}}{}
\@ifundefined{definition}{\newtheorem{definition}[theorem]{Definition}}{}
\@ifundefined{lemma}{\newtheorem{lemma}[theorem]{Lemma}}{}
\makeatother
\title[On Completion Times under Memoryless Catastrophe]{On Completion Times under Memoryless Catastrophe}
\author{Sichen Wang, Zhipeng Lu}
\date{Source: arXiv:2609.16566v1 (CC BY 4.0)}
\begin{document}
\maketitle

\noindent\emph{Source and licence.} On Completion Times under Memoryless Catastrophe. Version arXiv:2609.16566v1 (CC BY 4.0), distributed under the Creative Commons Attribution 4.0 International licence, as recorded in the arXiv licence metadata for this exact version and shown on the abstract page https://arxiv.org/abs/2609.16566v1. Full rights notice: licenses/arxiv-2609.16566-CC-BY-4.0.txt.\par\smallskip

\noindent\textit{Editorial scope.} Counted unit: the Theorem whose source label is \texttt{thm:ct-char} in the article (source0.tex lines 1153--1166, source label \texttt{thm:ct-char}), delivered together with the article's own complete original proof of it (source0.tex lines 1168--1278, one \texttt{proof} environment of 111 lines). No mathematics is written, repaired or re-derived: statement and proof are the article's own text line for line. Required context retained uncounted: def:affine (lines 320--328); lem:affine-ratio-rec (lines 403--425); def:hazard (lines 502--508); def:geom-tail (lines 548--554); thm:char (lines 596--604); eq:p-ct (lines 1057--1059); thm:ct-pgf (lines 1066--1071); def:ct-affine (lines 1134--1140). Every definition, display and label of those lines is retained unchanged. Omitted from this package and not redistributed: 1--319, 329--402, 426--501, 509--547, 555--595, 605--1056, 1060--1065, 1072--1133, 1141--1152, 1167--1167, 1279--1793. Nothing inside a retained excerpt is omitted or altered: every retained body is delivered exactly as the article's lines are written, so no omission receipt of any kind accompanies this package. Citations: the retained excerpts carry no citation command at all, so no bibliography entry of the article is transcribed and no reference is redistributed. Reference adaptation: none was needed and none was made; every reference inside the delivered unit resolves inside it, so no unresolved reference or citation is produced. Printed slips kept literal: no printed word, symbol, hypothesis or number of the counted statement or of its proof was altered. Layout and class: the article's own class is \texttt{article} with packages including fontenc, amsmath, amsthm, mathtools, libertine, newtxmath, microtype, geometry, hyperref, enumerate, booktabs, array, float, natbib; the wrapper is the lane's pinned \texttt{amsart} editorial wrapper at 10pt on A4, which restates the theorem-like environments the retained text needs and adds the article's own macro definitions that the retained text needs (none), with extra wrapper packages (bm, bbm, dsfont, xspace, cancel, mathtools, enumitem). The delivered numbering is the unit's own. Assets: no retained span contains a figure environment, a graphics-inclusion command, a PDF-page inclusion, a table or a TikZ picture; no image file is copied into this package, the package declares no asset dependency and no image file is redistributed with it. Licence: version arXiv:2609.16566v1 of this article is distributed under the Creative Commons Attribution 4.0 International licence, as recorded in the arXiv licence metadata for this exact version and shown on the abstract page https://arxiv.org/abs/2609.16566v1. A full rights notice is delivered at licenses/arxiv-2609.16566-CC-BY-4.0.txt. Characters: every retained excerpt is pure ASCII, so the delivered engine, XeLaTeX with its default Latin Modern text font, reports no missing character.
\begin{definition}[Uniform affine PGF structure]\label{def:affine}
A catastrophe mechanism has uniform affine PGF structure if there exist functions
\(a(w),b(w),c(w),d(w)\), analytic in a neighborhood of \(w=1\), depending only on the mechanism
(not on the base process \(T_0\)), and a linear function \(u(w)\) with \(u(1)\in(0,1)\), such that for all base processes~$T_0$,
\[
  G_T(w) = \frac{a(w)\,g(u(w)) + b(w)}{c(w)\,g(u(w)) + d(w)},
\]
with $a$ and $c$ not identically zero, and the non-degeneracy condition $c(1)\,g(u(1)) + d(1) \neq 0$ for all admissible~$g$.\label{rem:nondegen}\footnote{Non-degeneracy excludes pathological representations obtained by multiplying numerator and denominator by a factor vanishing at $w = 1$, which would create a $0/0$ indeterminacy and invalidate the coefficient-matching argument.}
\end{definition}

\begin{lemma}[Affine-ratio recursion]\label{lem:affine-ratio-rec}
Let
\[
  R(z)=\frac{a(z)h(u(z))+b(z)}{c(z)h(u(z))+d(z)},
  \qquad
  u(z)=u_0+\kappa(z-z_*),
\]
where $a,b,c,d$ are $C^\infty$ near $z_*$, and $a(z_*)=c(z_*), b(z_*)=d(z_*), D(z_*):=c(z_*)h(u_0)+d(z_*)\neq 0$.
Set
\[
  N(z)=a(z)h(u(z))+b(z),\qquad D(z)=c(z)h(u(z))+d(z).
\]
Then, for every $k\ge 1$, $\Delta_k:=N^{(k)}(z_*)-D^{(k)}(z_*)$ depends only on $h(u_0),h'(u_0),\dots,h^{(k-1)}(u_0)$, and
\[
  R^{(k)}(z_*)
  =
  \frac{1}{D(z_*)}
  \left(
    \Delta_k-\sum_{j=1}^{k-1}\binom{k}{j}D^{(j)}(z_*)R^{(k-j)}(z_*)
  \right).
\]
Consequently, $R^{(k)}(z_*)$ depends only on the $(k-1)$-jet of $h$ at $u_0$.
\end{lemma}

\begin{definition}[Age-based catastrophe mechanism]\label{def:hazard}
An \emph{age-based catastrophe mechanism} is specified by a hazard sequence $(h_t)_{t \geq 1}$ with $h_t \in (0,1)$ for all $t$.  Here $h_t$ is the conditional probability of catastrophe at step $t$, given that no catastrophe has occurred in steps $1, \ldots, t-1$.  The associated \emph{survival function} is
\begin{equation}\label{eq:survival}
  S(0) = 1, \qquad S(t) = \prod_{i=1}^{t}(1 - h_i), \quad t \geq 1.
\end{equation}
Thus $S(t)$ is the probability of surviving $t$ consecutive steps without catastrophe.  The requirement $h_t \in (0,1)$ ensures $S(t) > 0$ for all $t$ (every step has a positive probability of both survival and catastrophe).
\end{definition}

\begin{definition}[Geometric-tail catastrophe]\label{def:geom-tail}
A hazard sequence $(h_t)_{t \geq 1}$ has \emph{geometric tail} if there exist $a \in (0,1)$ and $q \in (0,1)$ such that
\[
  h_1 = 1-a, \qquad h_t = q \quad \text{for all } t \geq 2.
\]
Equivalently, $S(t) = a\,s^{t-1}$ for $t \geq 1$, where $s = 1-q$.  The parameter $a = S(1) = 1-h_1$ is the first-step survival probability.  The special case $a = s$ (i.e., $h_1 = q$) recovers memoryless catastrophe.
\end{definition}

\begin{theorem}[Characterization]\label{thm:char}
Among age-based catastrophe mechanisms (Definition~\ref{def:hazard}), the following are equivalent:
\begin{enumerate}[\rm (A)]
  \item \emph{Geometric-tail catastrophe}: $S(t) = a\,s^{t-1}$ for $t \geq 1$ (Definition~\ref{def:geom-tail}).
  \item \emph{Uniform affine PGF structure}: Definition~\ref{def:affine}.
  \item \emph{Derivative reduction}: there exists $z_0 \in (0,1)$ such that for every $k \geq 1$ and every finitely-supported $T_0$, the $k$-th factorial moment $\eta_k$ depends on $T_0$ only through $\{g^{(j)}(z_0)\}_{j=0}^{k-1}$.
  \item \emph{Proportional single-point sufficiency}: there exist $z_0 \in (0,1)$ and $\lambda > 0$ such that $p = \lambda\,g(z_0)$ for all finitely-supported $T_0$.
\end{enumerate}
\end{theorem}

\begin{equation}\label{eq:p-ct}
  p = P(T_0 < \tau) = E[e^{-rT_0}] = \hat{f}_0(r),
\end{equation}

\begin{theorem}[Laplace transform under Poisson resetting]\label{thm:ct-pgf}
For any base process $T_0$ under Poisson resetting at rate $r$,
\begin{equation}\label{eq:ct-lt}
  \hat{f}_T(\lambda) = \frac{(\lambda + r)\,\hat{f}_0(\lambda + r)}{\lambda + r\,\hat{f}_0(\lambda + r)}, \qquad \sigma := \lambda + r.
\end{equation}
\end{theorem}

\begin{definition}[Uniform affine Laplace structure]\label{def:ct-affine}
A continuous-time catastrophe mechanism has \emph{uniform affine Laplace structure} if there exist functions \(a(\lambda),b(\lambda),c(\lambda),d(\lambda)\), analytic in a neighborhood of \(\lambda=0\), depending only on the mechanism and a linear function $\sigma(\lambda) = \lambda + r$ ($r > 0$), such that for all base processes $T_0$,
\[
  \hat{f}_T(\lambda) = \frac{a(\lambda)\,\hat{f}_0(\sigma) + b(\lambda)}{c(\lambda)\,\hat{f}_0(\sigma) + d(\lambda)},
\]
with $a, c$ not identically zero and $c(0)\hat{f}_0(r) + d(0) \neq 0$ for all admissible $\hat{f}_0$.
\end{definition}

\begin{theorem}[Continuous-time characterization]\label{thm:ct-char}
Among continuous-time age-based catastrophe mechanisms specified by a measurable, locally integrable hazard rate $h: (0,\infty) \to (0,\infty)$ (so that $S(x) = \exp(-\int_0^x h(t)\,dt)$ is well-defined, strictly decreasing, and absolutely continuous), the following are equivalent:
\begin{enumerate}[\rm (a)]
  \item $S(x) = e^{-rx}$ for all $x > 0$ (Poisson resetting; equivalently,
  $h(t) = r$ for a.e.\ $t > 0$).
  \item Uniform affine Laplace structure (Definition~\ref{def:ct-affine}).
  \item Derivative reduction: there exists $r_0 > 0$ such that for every
  positive, almost surely finite $T_0$ and every $k \ge 1$, the $k$-th
  moment of $T$ is determined by
  $\{\hat f_0^{(j)}(r_0)\}_{j=0}^{k-1}$.
  \item Single-point sufficiency: there exists $r_0 > 0$ such that
  $p = \hat f_0(r_0)$ for every positive, almost surely finite $T_0$.
\end{enumerate}
\end{theorem}

\begin{proof}
(a)$\Rightarrow$(b): Theorem~\ref{thm:ct-pgf}.

(b)$\Rightarrow$(c): Suppose
\[
  \hat f_T(\lambda)=\frac{a(\lambda)\hat f_0(\sigma(\lambda))+b(\lambda)}
  {c(\lambda)\hat f_0(\sigma(\lambda))+d(\lambda)},
  \qquad \sigma(\lambda)=\lambda+r_0.
\]
Applying $\hat f_T(0)=1$ to point masses $T_0=\delta_x$ and letting $x\to\infty$ gives $b(0)=d(0)$; back-substitution yields $a(0)=c(0)$. By non-degeneracy, Lemma~\ref{lem:affine-ratio-rec} applies with $R=\hat f_T$, $h=\hat f_0$, $z_*=0$, and $u_0=r_0$. Hence $\hat f_T^{(k)}(0)$ depends only on $\{\hat f_0^{(j)}(r_0)\}_{j=0}^{k-1}$, which is exactly the derivative-reduction property.

(a)$\Rightarrow$(d): By definition~\eqref{eq:p-ct},
\[
  p=P(T_0<\tau)=E[e^{-rT_0}]=\hat f_0(r).
\]

(d)$\Rightarrow$(a): Under (d), for every positive, almost surely finite
$T_0$ we have
\[
  E[S(T_0)] = p = \hat f_0(r_0)=E[e^{-r_0 T_0}].
\]
Taking $T_0=\delta_x$ gives $S(x)=e^{-r_0 x}$ for every $x>0$.

(c)$\Rightarrow$(a): We use only the $k=1$ case.  Thus there exist
$r_0>0$ and a function $\psi:(0,1)\to\mathbb R$ such that
\[
  E[T]=\psi(\hat f_0(r_0))
\]
for every positive, almost surely finite $T_0$.
For a general continuous-time age-based catastrophe mechanism,
\[
  p=E[S(T_0)],\qquad
  E[\ell\,|\,T_0=x]=F(x):=\int_0^x S(t)\,dt,
\qquad
  E[T]=\frac{E[\ell]}{p}
      =\frac{E[F(T_0)]}{E[S(T_0)]}.
\]
Hence
\begin{equation}\label{eq:ct-char-psi}
  \frac{E[F(T_0)]}{E[S(T_0)]}
  =\psi\!\bigl(E[e^{-r_0 T_0}]\bigr)
\end{equation}
for all positive, almost surely finite $T_0$.

\emph{Step 1: point masses determine $\psi$ on $(0,1)$.}
Taking $T_0=\delta_x$ in~\eqref{eq:ct-char-psi} gives
\[
  \psi(e^{-r_0 x})=\frac{F(x)}{S(x)}, \qquad x>0.
\]
Since $x\mapsto e^{-r_0 x}$ is a bijection $(0,\infty)\to(0,1)$, point
masses determine $\psi$ on all of $(0,1)$.

\emph{Step 2: two-point laws force $\psi$ to be M\"obius.}
Take $T_0=\alpha\delta_a+(1-\alpha)\delta_b$ with $0<a<b$.  Writing
\[
  v=\alpha e^{-r_0 a}+(1-\alpha)e^{-r_0 b}\in[e^{-r_0 b},e^{-r_0 a}],
\]
equation~\eqref{eq:ct-char-psi} becomes
\[
  \psi(v)
  =
  \frac{\alpha F(a)+(1-\alpha)F(b)}
       {\alpha S(a)+(1-\alpha)S(b)}.
\]
Solving for $\alpha$ in terms of $v$ shows that the right-hand side is a
linear-fractional function of $v$, namely
\[
  \psi(v)
  =
  \frac{(F(a)-F(b))\,v + F(b)e^{-r_0 a}-F(a)e^{-r_0 b}}
       {(S(a)-S(b))\,v + S(b)e^{-r_0 a}-S(a)e^{-r_0 b}}.
\]
Thus $\psi$ is M\"obius on every interval
$[e^{-r_0 b},e^{-r_0 a}]$.

\emph{Step 3: overlapping intervals glue the local M\"obius maps globally.}
Exactly as in Steps~2--4 of Theorem~\ref{thm:char}, the overlapping-interval
consistency argument shows that all these local M\"obius maps coincide with a
single global one:
\[
  \psi(v)=\frac{Av+B}{Cv+D},\qquad v\in(0,1).
\]
Because $S$ is strictly decreasing, the denominator slope in the two-point
formula is nonzero, so the global denominator coefficient satisfies $C\neq0$.

\emph{Step 4: coefficient consistency forces $F=\rho S + c$.}
Comparing the coefficients of $v$ in the two-point representation with the
global M\"obius map shows that
\[
  \frac{F(a)-F(b)}{S(a)-S(b)}
\]
is independent of the pair $(a,b)$ with $a\neq b$.  Denoting this constant by
$\rho$, we obtain
\[
  F(x)=\rho\,S(x)+c,\qquad x>0
\]
for some constant $c$.

\emph{Step 5: differentiation yields constant hazard everywhere.}
Since $F(x)=\int_0^x S(t)\,dt$ is absolutely continuous and
$F'(x)=S(x)$ a.e., differentiating $F=\rho S+c$ gives
\[
  S(x)=\rho\,S'(x)\qquad\text{for a.e.\ }x>0.
\]
As $S'(x)=-h(x)S(x)$ a.e.\ and $S(x)>0$, we conclude
\[
  h(x)=-\frac1\rho=:r
  \qquad\text{for a.e.\ }x>0.
\]
Hence $S(x)=e^{-rx}$ for all $x>0$, which is exactly (a).
\end{proof}

\end{document}
